\documentclass[10pt,a4paper]{amsart}
\usepackage[margin=19mm]{geometry}
\usepackage{amsmath,amssymb,mathtools}
\usepackage[scr=rsfs,cal=euler]{mathalfa}
\usepackage{xfrac}
\usepackage[unicode,hidelinks,bookmarks=false]{hyperref}
\newcommand{\ie}{i.e.\ }
\newcommand{\cf}{cf.\ }
\newcommand{\resp}{respectively\ }
\newcommand{\Subsection}{\subsection}
\providecommand{\nobreakdash}{\nobreak\hbox{-}}
\setlength{\emergencystretch}{2em}
\multlinegap0pt
\usepackage{dsfont}
\usepackage{amssymb}
\newtheorem{lemma}{Lemma}
\newtheorem{proposition}{Proposition}
\newcommand{\gr}{\geqslant}
\newcommand{\ls}{\leqslant}
\newcommand{\vol}{\mathrm{vol}}
\title{\bf Regular functional covering numbers}
\author{Apostolos Giannopoulos, Natalia Tziotziou}
\date{Source: arXiv:2512.04301v2 (CC BY 4.0)}
\begin{document}
\maketitle

\noindent\emph{Source and licence.} \bf Regular functional covering numbers. Version arXiv:2512.04301v2 (CC BY 4.0), distributed by arXiv under the Creative Commons Attribution 4.0 International licence, http://creativecommons.org/licenses/by/4.0/, as recorded in the arXiv licence metadata for this exact version and shown on the abstract page https://arxiv.org/abs/2512.04301v2. Full licence notice: \texttt{licenses/arxiv-2512.04301-CC-BY-4.0.txt}.\par\smallskip

\noindent\textit{Editorial scope.} Counted unit: the article's proposition prop:fcover-1 (source0.tex lines 719--727). The article's complete original proof of it is delivered with it (source0.tex lines 729--791, one \texttt{proof} environment of 63 lines). No mathematics is written, repaired or re-derived: the statement and the proof are the article's own lines, in the article's own notation, and no retained line is substituted or re-flowed.  Retained uncounted as the source-local context the proof needs: lines 539--541; lines 557--566.  Nothing was omitted from any retained span, so the package carries no omission record.  No retained line contains a citation, so the wrapper carries no bibliography and no bibliography entry.  No retained line contains a figure environment, an image-inclusion command or drawing code, so no image is delivered and the package declares no asset dependency.  Editorial formatting only, in addition: an AMS \texttt{amsart} preamble that restates the article's own declarations and macros used by the retained lines, namely the environments \texttt{lemma}, \texttt{proposition} on separate counters, together with the standard packages those lines need.  The retained environments print in this document as Lemma 1, Proposition 1 in order of appearance, whereas the article prints the same items with the numbering of its own full document.  The complete native source of the article as distributed in the arXiv e-print for this exact version is bundled unchanged as \texttt{sources/190324/source0.tex}.
\begin{equation} \label{eq:indicator-covering}
N(\mathds{1}_K,\mathds{1}_T) \ls N(K,T).
\end{equation}

\begin{lemma}\label{lem:properties}
Let $f,g,h,f_i,g_i \in LC_g(\mathbb{R}^n)$. Then, for any $a,b>0$ and $T\in GL_n$,
\begin{enumerate}
\item[{\rm (i)}] $N(af,bg)=\dfrac{a}{b}N(f,g)$.
\item[{\rm (ii)}] $N(f\circ T,\, g\circ T)=N(f,g)$.
\item[{\rm (iii)}] $N(f_1+f_2,g)\ls N(f_1,g)+N(f_2,g)$.
\item[{\rm (iv)}] If $f_1 \ls f_2$ and $g_1 \gr g_2$, then $N(f_1,g_1)\ls N(f_2,g_2)$.
\item[{\rm (v)}] $N(f,g)\ls N(f,h)\,N(h,g)$.
\end{enumerate}
\end{lemma}

\begin{proposition}\label{prop:fcover-1}
Let $f:\mathbb{R}^n\to[0,\infty)$ be a geometric log-concave function.
If $\vol_n(R_f)^{1/n}\approx 1$ and $R_f$ satisfies 
\begin{equation}\max\Bigl\{N(R_f,\, t r_f B_2^n),\,
N(B_2^n,\, t r_f (R_f)^{\circ})\Bigr\}\ls \exp\left( \tfrac{c_1 \kappa_n^2 n}{t^2}\right),
\label{eq:R-covering-a}\end{equation}
for some $\kappa_n\gr 1$ and every $t\gr 1$, then
$$N(f,t\odot g)\ls \exp\!\left(\frac{\kappa_n n}{t}\right)\qquad\text{for all } t\gr 1.$$
\end{proposition}

\begin{proof}
We first claim that
\begin{equation}\label{eq:prop-fcover-1.1}
f(x)\ls \mathds{1}_{R_f}(x)+\exp\!\left(-50n\,\|x\|_{R_f}\right),\qquad x\in\mathbb{R}^n.
\end{equation}
If $x\in R_f$, then $f(x)\ls 1=\mathds{1}_{R_f}(x)$. If $x\notin R_f$, then $\lambda:=\|x\|_{R_f}>1$ and by definition
$f(x/\lambda)=\exp(-50n)f(0)$. Log-concavity yields
$$f(x/\lambda)\gr f(x)^{\frac{1}{\lambda}}f(0)^{1-\frac{1}{\lambda}}$$
and hence
\begin{equation}\label{eq:no-need-1}f(x)\ls \left(\frac{f(x/\lambda)}{f(0)}\right)^{\lambda}f(0)=\exp\!\left(-50n\,\|x\|_{R_f}\right)f(0)
\ls \exp\!\left(-50n\,\|x\|_{R_f}\right).\end{equation}
Next, decompose $\exp(-50n\|x\|_{R_f})$ as
\begin{equation}\label{eq:fcover-exp}
\exp(-50n\|x\|_{R_f})=\sum_{k=0}^{\infty}e^{-50n\|x\|_{R_f}}
\,\mathds{1}_{\{k\ls 50n\|x\|_{R_f}<k+1\}}(x)
\ls \sum_{k=0}^{\infty}e^{-k}\,\mathds{1}_{\frac{k+1}{50n}R_f}(x).
\end{equation}
Combining \eqref{eq:prop-fcover-1.1}, \eqref{eq:fcover-exp}, and Lemma~\ref{lem:properties}\,(iii),
\begin{equation}\label{eq:fcover-sum}
N(f,t\odot g)\ls N(\mathds{1}_{R_f},t\odot g)+\sum_{k=0}^{\infty}e^{-k}\,
N(\mathds{1}_{\frac{k+1}{50n}R_f},\,t\odot g).
\end{equation}

We begin with $N(\mathds{1}_{R_f},t\odot g)$.
By submultiplicativity,
\begin{equation}\label{eq:fcover-first}
N(\mathds{1}_{R_f},t\odot g)\ls N(\mathds{1}_{R_f},\,\mathds{1}_{\sqrt{\kappa_nt}r_f B_2^n})
\,N(\mathds{1}_{\sqrt{\kappa_nt}r_f B_2^n},\,t\odot g).
\end{equation}
By \eqref{eq:indicator-covering} and \eqref{eq:R-covering-a},
\begin{equation}\label{eq:fcover-first-1}
N(\mathds{1}_{R_f},\,\mathds{1}_{\sqrt{\kappa_nt}r_f B_2^n})\ls \exp\!\left(\frac{c\kappa_n n}{t}\right).
\end{equation}
Since $\vol(R_f)^{1/n}\approx 1$, we have $r_f\approx\sqrt{n}$, and for $x\in \sqrt{\kappa_nt}r_f B_2^n$,
$$\exp\!\left(-\frac{|x|^2}{2t^2}\right)\gr\exp\left(-\frac{\kappa_nr_f^2}{2t}\right) 
\gr\exp(-c\kappa_nn/t),$$
where $c>0$ is an absolute constant. Thus
$$\mathds{1}_{\sqrt{\kappa_nt}r_f B_2^n}(x)\ls e^{c\kappa_nn/t}(t\odot g)(x),$$
and hence (by Lemma~\ref{lem:properties}\,(i))
\begin{equation}\label{eq:fcover-first-2}
N(\mathds{1}_{\sqrt{\kappa_nt}r_f B_2^n},t\odot g)\ls e^{c\kappa_nn/t}.
\end{equation}
Combining \eqref{eq:fcover-first-1}--\eqref{eq:fcover-first-2} gives
\begin{equation}\label{eq:fcover-first-final}
N(\mathds{1}_{R_f},t\odot g)\ls \exp\!\left(\frac{c_1\kappa_n n}{t}\right).
\end{equation}

Next, we give an upper bound for the sum in \eqref{eq:fcover-sum}. For $k\gr 0$, using \eqref{eq:fcover-first-2}
and submultiplicativity, we write
\begin{align*}N(\mathds{1}_{\frac{k+1}{50n}R_f},t\odot g) &\ls N(\mathds{1}_{\frac{k+1}{50n}R_f},\mathds{1}_{\sqrt{\kappa_nt}r_f B_2^n})
\,N(\mathds{1}_{\sqrt{\kappa_nt}r_f B_2^n},t\odot g)\\
&\ls e^{c\kappa_nn/t} N(R_f,\sqrt{\kappa_nt}r_f B_2^n)\,N\!\left(B_2^n,\tfrac{50n}{k+1}B_2^n\right)\,.
\end{align*}
Using \eqref{eq:R-covering-a} and the bound $N(B_2^n,\lambda B_2^n)\ls (1+2/\lambda)^n$,
$$N(\mathds{1}_{\frac{k+1}{50n}R_f},t\odot g)\ls \exp\!\left(\frac{c_2\kappa_n n}{t}\right)e^{\frac{k+1}{25}}.$$
It follows that
$$\sum_{k=0}^\infty e^{-k}N(\mathds{1}_{\frac{k+1}{50n}R_f},t\odot g)\ls 
C_1\exp\!\left(\frac{c_2\kappa_nn}{t}\right),$$
and together with \eqref{eq:fcover-sum} and \eqref{eq:fcover-first-final},
$$N(f,t\odot g)\ls \exp\!\left(\frac{c_1\kappa_n n}{t}\right)
+C_1\exp\!\left(\frac{c_2\kappa_n n}{t}\right),$$
which gives the required bound.
\end{proof}

\end{document}
